\documentclass[11pt,a4paper]{article}
\usepackage[utf8]{inputenc}
\usepackage{amsmath,amssymb,amsthm}
\usepackage{algorithm}
\usepackage{algpseudocode}
\usepackage{booktabs}
\usepackage{hyperref}
\usepackage{geometry}
\geometry{margin=1in}

\title{Deformable Convolutional Networks: Non-Rigid Spatial Sampling, Bilinear Interpolation, and Modulation Mechanics}
\author{Automorphic Intelligence Architecture Specifications}
\date{\today}

\newtheorem{theorem}{Theorem}
\newtheorem{lemma}[theorem]{Lemma}
\newtheorem{proposition}[theorem]{Proposition}
\theoremstyle{definition}
\newtheorem{definition}{Definition}
\newtheorem{remark}{Remark}

\begin{document}

\maketitle

\begin{abstract}
Standard convolutional operations operate over fixed, rigid discrete spatial grids, constraining their receptive fields to uniform rectangular geometries irrespective of non-rigid object deformations or scale variations. Deformable Convolution (DCN) generalizes 2D convolution by augmenting regular spatial sampling grids with learned 2D continuous offsets and amplitude modulation factors. This document formalizes the continuous spatial interpolation operator, derives sub-gradient propagation equations through bilinear kernels, establishes computational bounds, presents complete pseudocode, and walks through a verified numerical example.
\end{abstract}

\section{Notation and Mathematical Preliminaries}

Let $x \in \mathbb{R}^{C_{\text{in}} \times H \times W}$ denote the input feature map and let $\mathcal{R}$ denote the regular spatial grid receptive field (e.g., $3 \times 3$ kernel grid $\mathcal{R} = \{(-1,-1), (-1,0), \dots, (1,1)\}$ with $|\mathcal{R}| = K = 9$).

\begin{table}[h]
\centering
\caption{Mathematical Notation for Deformable Convolution}
\begin{tabular}{lll}
\toprule
\textbf{Symbol} & \textbf{Tensor Shape} & \textbf{Semantic Description} \\
\midrule
$p_0$ & $\mathbb{Z}^2$ & Target spatial evaluation coordinate $(h, w)$ \\
$p_k$ & $\mathbb{Z}^2$ & Nominal kernel grid offset $\in \mathcal{R}$ \\
$\Delta p_k(p_0)$ & $\mathbb{R}^2$ & Learned continuous spatial displacement $(\Delta y_k, \Delta x_k)$ \\
$m_k(p_0)$ & $[0, 1]$ & Learned scalar modulation amplitude (DCNv2) \\
$w_k$ & $\mathbb{R}^{C_{\text{out}} \times C_{\text{in}}}$ & Filter weight matrix for kernel index $k$ \\
$G(q, p)$ & $\mathbb{R}$ & Bilinear spatial interpolation kernel at point $q$ for target $p$ \\
$y(p_0)$ & $\mathbb{R}^{C_{\text{out}}}$ & Output feature vector at spatial location $p_0$ \\
\bottomrule
\end{tabular}
\end{table}

\section{Problem Definition and Core Concepts}

\subsection{3-Level Explanation}
\begin{enumerate}
    \item \textbf{Intuitive (High School level):} A normal camera lens looks through a rigid $3 \times 3$ square of pixels. If an object is curved or tilted (like a cat stretching or a snake curling), rigid squares miss the object. Deformable convolution lets each of the 9 points bend, stretch, and move directly onto the target object.
    \item \textbf{Engineering Level:} DCN introduces a small auxiliary convolutional layer that predicts $2K$ spatial offsets $\Delta p_k$ (and optionally $K$ modulation weights $m_k \in [0, 1]$ via sigmoid). Features are sampled at fractional coordinates $(p_0 + p_k + \Delta p_k)$ using bilinear interpolation kernels before multiplying by the main convolutional weights.
    \item \textbf{Mathematical Level:} For spatial location $p_0$, the modulated deformable convolution operator is defined by:
    \begin{equation}
    y(p_0) = \sum_{k=1}^K w_k \cdot m_k(p_0) \cdot x(p_0 + p_k + \Delta p_k(p_0))
    \end{equation}
    where fractional sampling is evaluated via continuous 2D bilinear basis functions:
    \begin{equation}
    x(p) = \sum_{q \in \mathbb{Z}^2} \max(0, 1 - |q_x - p_x|) \cdot \max(0, 1 - |q_y - p_y|) \cdot x(q)
    \end{equation}
\end{enumerate}

\subsection{5-Step Algorithmic Framework}
\begin{enumerate}
    \item \textbf{Offset and Modulation Prediction:} Compute $\Delta p = \text{Conv}(x; W_{\text{off}})$ and $m = \sigma(\text{Conv}(x; W_{\text{mod}}))$.
    \item \textbf{Fractional Sampling Point Computation:} For each kernel point $k \in \{1,\dots, K\}$, compute continuous coordinates $p = p_0 + p_k + \Delta p_k(p_0)$.
    \item \textbf{Bilinear Neighborhood Gathering:} Identify the four integer grid neighbors $\{q_{11}, q_{12}, q_{21}, q_{22}\}$ surrounding $p$.
    \item \textbf{Interpolated Value Assembly:} Compute weighted sum of neighboring pixel intensities using distance weights $(1 - |q_x - p_x|)(1 - |q_y - p_y|)$.
    \item \textbf{Modulated Kernel Contraction:} Multiply interpolated values by modulation $m_k$ and kernel weights $w_k$, accumulating across input channels.
\end{enumerate}

\section{Algorithm Specification}

\begin{algorithm}[H]
\caption{Modulated Deformable Convolution Forward Pass}
\label{alg:deform_conv}
\begin{algorithmic}[1]
\Require Input tensor $x \in \mathbb{R}^{C_{\text{in}} \times H \times W}$, learned offsets $\Delta p \in \mathbb{R}^{K \times 2 \times H \times W}$, modulations $m \in \mathbb{R}^{K \times H \times W}$, filter weights $W \in \mathbb{R}^{C_{\text{out}} \times C_{\text{in}} \times 3 \times 3}$, bias $b \in \mathbb{R}^{C_{\text{out}}}$.
\Ensure Output tensor $y \in \mathbb{R}^{C_{\text{out}} \times H \times W}$.
\State $\mathcal{R} \gets \{(-1,-1), (-1,0), (-1,1), (0,-1), (0,0), (0,1), (1,-1), (1,0), (1,1)\}$
\For{$c_o = 1$ \textbf{to} $C_{\text{out}}$}
    \For{$i = 1$ \textbf{to} $H$}
        \For{$j = 1$ \textbf{to} $W$}
            \State $v \gets b[c_o]$
            \For{$k = 1$ \textbf{to} $K$}
                \State $(g_y, g_x) \gets \mathcal{R}[k]$
                \State $p_y \gets (i + g_y) + \Delta p[k, 0, i, j]$
                \State $p_x \gets (j + g_x) + \Delta p[k, 1, i, j]$
                \State $m_{\text{val}} \gets m[k, i, j]$
                \For{$c_i = 1$ \textbf{to} $C_{\text{in}}$}
                    \State $x_{\text{sampled}} \gets \text{BilinearSample2D}(x[c_i], p_y, p_x)$
                    \State $v \gets v + W[c_o, c_i, g_y+1, g_x+1] \cdot m_{\text{val}} \cdot x_{\text{sampled}}$
                \EndFor
            \EndFor
            \State $y[c_o, i, j] \gets v$
        \EndFor
    \EndFor
\EndFor
\State \Return $y$
\end{algorithmic}
\end{algorithm}

\section{Theoretical Guarantees and Differentiability}

\begin{theorem}[Differentiability of Bilinear Offset Propagation]
\label{thm:differentiability}
The loss gradient with respect to the spatial offset $\Delta p_x$ is continuous almost everywhere (piecewise linear in spatial coordinates) and is given analytically by:
\begin{equation}
\frac{\partial y(p_0)}{\partial \Delta p_{k,x}} = \sum_{c_i} w_{k, c_i} m_k(p_0) \sum_{q} \frac{\partial G(q, p_0 + p_k + \Delta p_k)}{\partial p_x} x(q)
\end{equation}
where:
\begin{equation}
\frac{\partial G(q, p)}{\partial p_x} = \begin{cases}
\text{sgn}(q_x - p_x) \max(0, 1 - |q_y - p_y|) & \text{if } |q_x - p_x| < 1 \text{ and } q_x \ne p_x \\
0 & \text{otherwise}
\end{cases}
\end{equation}
\end{theorem}

\begin{proof}
Directly follows by differentiating the bilinear kernel $G(q, p) = g(q_x, p_x) g(q_y, p_y)$ with respect to $p_x$, where $g(u, v) = \max(0, 1 - |u - v|)$.
\end{proof}

\section{Complexity Analysis}

\begin{itemize}
    \item \textbf{Time Complexity:}
    \begin{itemize}
        \item Bilinear sampling: 4 memory reads + 6 FLOPs per sampling point $\to \mathcal{O}(4 \cdot K \cdot C_{\text{in}} \cdot H \cdot W)$.
        \item Kernel multiplication and modulation: $\mathcal{O}(K \cdot C_{\text{in}} \cdot C_{\text{out}} \cdot H \cdot W)$.
        \item Total Time: $\mathcal{O}\left( K \cdot C_{\text{in}} \cdot C_{\text{out}} \cdot H \cdot W \right)$ (equal FLOP count to standard conv, but memory-bound due to non-contiguous spatial gathers).
    \end{itemize}
    \item \textbf{Space Complexity:} $\mathcal{O}(3K \cdot H \cdot W)$ for offset and modulation field tensors.
\end{itemize}

\section{Exactness and Error Analysis}

The bilinear interpolation introduces first-order spatial interpolation error bounded by the second derivative of the feature map:
\begin{equation}
|x(p) - x_{\text{exact}}(p)| \le \frac{1}{8} \left( \|\partial_{xx}^2 x\|_{\infty} + \|\partial_{yy}^2 x\|_{\infty} \right)
\end{equation}

\section{Worked Numerical Example}

Let $C_{\text{in}}=1, C_{\text{out}}=1, H=2, W=2$, kernel size $1 \times 1$ ($K=1, p_1 = (0,0)$), weight $w = 1.0$, modulation $m = 1.0$.
\begin{equation}
x = \begin{bmatrix} 10.0 & 20.0 \\ 30.0 & 40.0 \end{bmatrix}
\end{equation}
Evaluate at coordinate $p_0 = (0, 0)$ with learned offset $\Delta p = (0.5, 0.5)$.
Target continuous coordinate $p = (0.5, 0.5)$:
\begin{itemize}
    \item $y_{\text{low}}=0, y_{\text{high}}=1 \implies l_y = 0.5, h_y = 0.5$.
    \item $x_{\text{low}}=0, x_{\text{high}}=1 \implies l_x = 0.5, h_x = 0.5$.
    \item Weights: $w_1 = 0.25, w_2 = 0.25, w_3 = 0.25, w_4 = 0.25$.
    \item $x(0.5, 0.5) = 0.25(10.0) + 0.25(20.0) + 0.25(30.0) + 0.25(40.0) = 2.5 + 5.0 + 7.5 + 10.0 = 25.0$.
    \item $y(0, 0) = 1.0 \times 1.0 \times 25.0 = 25.0$.
\end{itemize}

\section{Numerical Considerations and Edge Cases}

\begin{itemize}
    \item \textbf{Out-of-Bounds Offsets:} If continuous coordinate $p$ leaves the image boundaries $[0, H-1] \times [0, W-1]$, zero-padding boundary conditions apply smoothly ($x(p) = 0$).
    \item \textbf{Gradient Clipping on Offsets:} Unbounded offsets can cause sampling points to explode outside the feature map. Softsign or hyperbolic tangent scaling is recommended during training.
\end{itemize}

\section{References}
\begin{enumerate}
    \item Dai, J., Qi, H., Xiong, Y., Li, Y., Zhang, G., Hu, H., & Wei, Y. (2017). Deformable convolutional networks. \emph{ICCV}, 764--773.
    \item Zhu, X., Hu, H., Lin, S., & Dai, J. (2019). Deformable ConvNets v2: More deformable, better results. \emph{CVPR}, 9308--9316.
\end{enumerate}

\end{document}
